\section{Introduction}
\label{sec:intro}

Large language models (LLMs) classify text from a few labelled examples placed in the prompt, without any weight update \citep{brown2020language}.
Which examples are shown, and in what order, can move accuracy by as much as a change in model size \citep{liu2022makes,lu2022fantastically,zhao2021calibrate}, and with many classes the prompt must also carry the label space.

Conformal In-Context Learning (CICLe) was proposed by \citet{randl2024cicle} to make this decision smaller before the LLM sees it.
A small classifier is calibrated with class-conditional conformal prediction so that, for each input, it returns a set of candidate labels that contains the true label with a chosen probability.
The LLM then chooses among these candidates only, with examples retrieved from the candidate classes, and is not called when the set has one element.

CICLe was evaluated on one food-safety dataset.
A pilot study on four further datasets, run with the original CICLe prompts, reported that its benefit grows with class imbalance and that it helps some LLMs and hurts others.
The pilot protocol let only the zero-shot prompt list the admissible labels, drew few-shot and CICLe examples from pools of different size, and left the conformal sets unseeded, so that runs could not be paired.
The first flaw alone made up to 65\% of the outputs of some models unparseable, and the parse failures, not the method, decided which models appeared to benefit.
% ANONYMITY: the pilot is arXiv:2512.05732 by the same authors. Do not cite it. Keep the wording "a pilot study".
This paper replaces the pilot with a corrected protocol and asks a sharper question, namely how a small classifier should narrow the label space for an LLM, and when the LLM is worth calling at all.
We evaluate CICLe on five English datasets, six open LLMs with 3B to 8B parameters and two with 12B and 32B, with three data seeds and paired bootstrap tests over test instances throughout.
Our contributions are the following.
\begin{enumerate}
  \setlength\itemsep{0pt}
  \item \textbf{A corrected protocol.} Every prompt lists the admissible labels, every method retrieves from the same pool, and calibration splits, conformal sets and test samples are seeded. Rerunning the pilot's protocol shows that these choices change conclusions (Section~\ref{sec:res-protocol}).
  \item \textbf{A measured gain.} With wording, pool and seeds held equal, CICLe beats few-shot prompting in nine of ten dataset and retrieval cells by 0.5 to 1.8 macro-F1 points, and shortens Per-Class prompts by 11 to 45\% (Section~\ref{sec:res-gain}).
  \item \textbf{How to narrow.} At equal mean set size, class-conditional conformal sets match top-$k$, probability-mass and marginal conformal sets on balanced data. When the labelled pool is long-tailed, the alternatives lose coverage on the rare classes and 3 to 13 points, which the LLM cannot recover. The gain over few-shot prompting is largest, up to 6.5 points, when label names are nonsense words (Sections~\ref{sec:res-narrowing} and~\ref{sec:res-relabel}).
  \item \textbf{When to call the LLM.} With 1,600 balanced labels, a fine-tuned RoBERTa-base beats every 3B to 8B pipeline on four of five datasets, and on the fifth the base classifier alone does. Under a long-tailed or small pool the order flips. We give a decision rule and a rule for the miscoverage level (Sections~\ref{sec:res-baselines} and~\ref{sec:res-alpha}).
\end{enumerate}
We retract the pilot's imbalance claim. Imbalance widens the gap between narrowing rules, not between CICLe and few-shot prompting.
Code, prompts and per-instance records will be released.
% Anonymised repository placeholder: \url{https://anonymous.4open.science/r/XXXX}
